LSZ 约化把时序关联函数转化为实验散射振幅以后，同一个解析振幅还包含比单一反应道更多的信息。改变某条外腿的入射/出射解释会把散射、湮灭和对产生联系起来；真实中间态打开后出现的虚部，则通过复解析性约束阈值以下的实部。前者由交叉关系描述，后者由色散关系描述。

第一个问题的答案来自外腿的解析延拓。LSZ 振幅取自同一个截肢关联函数，只是不同物理过程把不同外腿解释为入射粒子、出射粒子或相应反粒子。改变这种解释时，四动量的符号和外线波函数必须一起改变。这一关系称为\emph{交叉}。

第二个问题来自复变量解析性。虚中间态尚不能处在真实质量壳上时，振幅可以平滑地随运动学变量变化；一旦某个真实多粒子通道打开，振幅在相应阈值处出现分支点，实轴两侧的边界值不再相同。这个差值就是割线不连续度，它的虚部由真实在壳中间态决定。若振幅在割线之外解析，Cauchy 公式便把这些边界值与其它运动学区域联系起来；为了控制高能端的积分，必要时要做有限次减法。这就是色散关系的来源。


交叉关系最清楚的记号是把四点截肢连通核统一写成“所有四动量均流入顶角”的形式
\begin{equation}
 \mathcal A(p_1,p_2,p_3,p_4),
 \qquad
 p_1+p_2+p_3+p_4=0.
 \label{eq:all-incoming-amplitude-dp9}
\end{equation}
对物理过程 $1+2\to3+4$，$p_1,p_2,p_3,p_4$ 都取正能物理四动量；出射外腿在全入射核中必须反向，因此
\begin{equation}
 \boxed{
 \mathcal M_{12\to34}
 =\mathcal A(p_1,p_2,-p_3,-p_4).}
 \label{eq:physical-from-all-incoming-dp9}
\end{equation}
若把某条外腿改解释为交叉过程中的反粒子，不能只机械地把 Mandelstam 变量换名；必须同时解析延拓该外腿四动量，并把 Dirac 外线的 $u$、$v$ 或矢量场的偏振波函数换成新物理态对应的外线因子。这个联合替换才保持同一截肢核与不同物理外态之间的对应。

\begin{note}[交叉关系的使用范围]
在逐阶微扰展开中，外腿延拓可直接作用于 Feynman 图的解析表达式。完全非微扰的交叉定理还依赖局域性、谱条件和解析延拓的全局性质。\eqnrange{eq:all-incoming-amplitude-dp9}{eq:physical-from-all-incoming-dp9} 所用的是标准局域量子场论中的外腿交叉规则，而不是对任意非局域模型的无条件恒等式。
\end{note}

转向第二个问题。固定其它运动学变量，把待研究的振幅或形状因子记作复变量 $z$ 的函数 $F(z)$。这里 $z$ 可取某个 Mandelstam 不变量，因而具有动量平方量纲。真实中间态通常从某个右手阈值 $s_R$ 开始产生割线；交叉道可以在负实轴方向产生左手割线，其最高端点记为 $s_L$。在两条割线之间选一个解析点 $z_0$ 作为减法点。

减法的目的很具体：如果 $F(z)$ 在无穷远衰减得不够快，就不能直接丢掉 Cauchy 大圆。取整数 $N\ge1$，用 $F$ 在 $z_0$ 的前 $N$ 阶 Taylor 数据构造
\begin{equation}
 \boxed{
 H_N(\zeta;z)
 \equiv
 \frac{
 F(\zeta)-\displaystyle\sum_{r=0}^{N-1}
 \frac{F^{(r)}(z_0)}{r!}(\zeta-z_0)^r
 }
 {(\zeta-z_0)^N(\zeta-z)}.}
 \label{eq:subtracted-cauchy-kernel-dp68}
\end{equation}
只要 $N$ 取得足够大，使 $H_N$ 在无穷远大圆上的积分消失，就可以把闭合轮廓收缩到实轴割线和孤立极点。于是得到 $N$ 次减法色散关系
\begin{equation}
 \boxed{
 \begin{aligned}
 F(z)
 &=\sum_{r=0}^{N-1}
 \frac{F^{(r)}(z_0)}{r!}(z-z_0)^r
 +\frac{(z-z_0)^N}{\pi}
 \int_{s_R}^{\infty}\!\dd s'\,
 \frac{\operatorname{Im}F(s'+\ii0^+)}{(s'-z_0)^N(s'-z)}
 \\
 &\quad
 +\frac{(z-z_0)^N}{\pi}
 \int_{-\infty}^{s_L}\!\dd s'\,
 \frac{\operatorname{Im}F(s'+\ii0^+)}{(s'-z_0)^N(s'-z)}
 +\text{孤立极点项}.
 \end{aligned}}
 \label{eq:general-subtracted-dispersion-dp9}
\end{equation}
这里的减法多项式由 $F(z_0),F'(z_0),\ldots$ 这样的低能输入或重整化条件确定；割线积分则来自真实中间态。二者的物理来源不同，不能把减法常数误认为由幺正性自动决定。

若只有从 $s_0$ 开始的右手割线，并取一次减法 $N=1$、$z_0=0$ 且 $F(0)=0$，\eqnrefs{eq:general-subtracted-dispersion-dp9} 退化为
\begin{equation}
 \boxed{
 F(z)
 =\frac{z}{\pi}
 \int_{s_0}^{\infty}\dd s'\,
 \frac{\operatorname{Im}F(s'+\ii0^+)}{s'(s'-z)}.}
 \label{eq:once-subtracted-single-cut-dp9}
\end{equation}
例如在真空极化中，时间样区域的真实 $e^-e^+$ 产生决定 $\operatorname{Im}F$，同一色散积分据此重建类空区域的实部。

对某些由光学定理保证谱权重 $\rho(s)\equiv\operatorname{Im}F(s+\ii0^+)\ge0$ 的标量投影，\eqnrefs{eq:once-subtracted-single-cut-dp9} 还给出一个直接的符号检查。令 $z=-Q^2<0$，只要相应谱矩收敛，就有
\begin{equation}
 \boxed{
 (-1)^n\frac{\dd^n}{\dd(Q^2)^n}F(-Q^2)\ge0,
 \qquad n\ge1.}
 \label{eq:spacelike-spectral-positivity-dp68}
\end{equation}
这不是任意振幅都满足的额外公理；它只适用于当前投影的谱密度确实非负的情形。

交叉关系与色散关系连接散射振幅的两个不同层面：前者把不同外态过程识别为同一截肢关联函数在不同质量壳边界上的取值，后者用割线虚部和必要的减法常数重建解析域中的其余部分。具体理论只需要把相应粒子质量、耦合与传播子代入这些一般结构。


\begin{derivation}[从\eqnrefs{eq:all-incoming-amplitude-dp9}到\eqnrefs{eq:physical-from-all-incoming-dp9}]
全入射约定把每一条外腿的四动量都定义成指向相互作用区。设全入射核的四个参数为 $k_i$，则正文\eqnrefs{eq:all-incoming-amplitude-dp9} 要求
\[
k_1+k_2+k_3+k_4=0.
\]
对物理过程 $1+2\to3+4$，实验四动量满足
\[
p_1+p_2=p_3+p_4.
\]
前两条腿本来就流入，所以
\[
k_1=p_1,\qquad k_2=p_2.
\]
末态 3、4 的物理箭头流出相互作用区；若仍要求 $k_i$ 全部向内，必须反号：
\[
k_3=-p_3,\qquad k_4=-p_4.
\]
于是全入射的动量守恒自动变成
\begin{align*}
k_1+k_2+k_3+k_4
&=p_1+p_2-p_3-p_4\\
&=0,
\end{align*}
与物理守恒完全一致。将四个 $k_i$ 代回同一个截肢四点核，便得到正文\eqnrefs{eq:physical-from-all-incoming-dp9}：
\[
\mathcal M_{12\to34}
=\mathcal A(p_1,p_2,-p_3,-p_4).
\]
\end{derivation}

\begin{derivation}[从\eqnrefs{eq:subtracted-cauchy-kernel-dp68}到\eqnrefs{eq:general-subtracted-dispersion-dp9}]
把轮廓收缩到两条实轴割线。实解析性给出

\begin{equation*}
 F(s'-\ii0^+)=F(s'+\ii0^+)^*,
\end{equation*}
所以割线两侧的差为
\begin{equation*}
 \operatorname{Disc}F(s')
 \equiv F(s'+\ii0^+)-F(s'-\ii0^+)
 =2\ii\operatorname{Im}F(s'+\ii0^+).
\end{equation*}
右手割线因此贡献
\begin{equation*}
 \frac{(z-z_0)^N}{\pi}
 \int_{s_R}^{\infty}\!\dd s'\,
 \frac{\operatorname{Im}F(s'+\ii0^+)}
 {(s'-z_0)^N(s'-z)},
\end{equation*}
左手割线同理给出从 $-\infty$ 到 $s_L$ 的积分。最后把预先减去的 Taylor 多项式加回，并把轮廓内未并入割线的孤立极点留作显式留数项，就得到正文\eqnrefs{eq:general-subtracted-dispersion-dp9}。

若 $z$ 逼近割线，分母 $1/(s'-z)$ 按 $z\to z+\ii0^+$ 理解。其主值部分给出实部，跨越极点产生的 $\ii\pi$ 项恢复原来的虚部，因此色散式同时编码了割线两侧的完整边界值。
\end{derivation}



\begin{derivation}[从\eqnrefs{eq:once-subtracted-single-cut-dp9}到\eqnrefs{eq:spacelike-spectral-positivity-dp68}]
这一结论只讨论谱密度 $\rho(s)\equiv\operatorname{Im}F(s+\ii0^+)\ge0$ 的投影。令 $z=-Q^2$、$Q^2>0$，正文\eqnrefs{eq:once-subtracted-single-cut-dp9} 变为
\begin{equation*}
F(-Q^2)
=-\frac{Q^2}{\pi}
\int_{s_0}^{\infty}\dd s'\,
\frac{\rho(s')}{s'(s'+Q^2)}.
\end{equation*}
利用恒等式
\begin{equation*}
\frac{Q^2}{s'(s'+Q^2)}
=\frac1{s'}-\frac1{s'+Q^2},
\end{equation*}
可以把所有 $Q^2$ 依赖集中到第二项：
\begin{equation*}
F(-Q^2)
=-\frac1\pi\int_{s_0}^{\infty}\frac{\rho(s')}{s'}\dd s'
+\frac1\pi\int_{s_0}^{\infty}\frac{\rho(s')}{s'+Q^2}\dd s'.
\end{equation*}
若相应谱矩收敛，可以逐次把 $Q^2$ 导数移入积分号。一次和二次导数分别为
\begin{align*}
\frac{\dd}{\dd Q^2}F(-Q^2)
&=-\frac1\pi\int_{s_0}^{\infty}\dd s'\,
\frac{\rho(s')}{(s'+Q^2)^2}\le0,\\
\frac{\dd^2}{\dd(Q^2)^2}F(-Q^2)
&=\frac{2}{\pi}\int_{s_0}^{\infty}\dd s'\,
\frac{\rho(s')}{(s'+Q^2)^3}\ge0.
\end{align*}
继续求导时，每次只会再产生一个负号和一个正的整数系数，因此归纳得到正文\eqnrefs{eq:spacelike-spectral-positivity-dp68}。它给出了一个很实用的自洽检查：若显式圈积分声称对应同一非负谱投影，却违反这一导数符号层级，就说明符号、减法或解析延拓中至少有一处不一致。
\end{derivation}

